\chapter*{Abstract}

This thesis presents \textit{MuseIQ}, an intelligent guide designed to improve the on-site visitor experience in museums and cultural spaces through low-cost and widely available technologies. The problem addressed is that many museums still rely on static or weakly adaptive mediation mechanisms, which limits personalization, accessibility, and the delivery of contextual information during the visit. In response, the main objective of this research is to propose a feasible technological solution capable of identifying the probable room or observation zone of the visitor, activating relevant content, and enabling conversational interaction supported by institutional curatorial knowledge.

The methodology follows an applied research approach focused on solution design. The study examines the Peruvian context of technology adoption and museum demand, reviews the state of the art in smart museums, indoor positioning, voice-based accessibility, and retrieval-augmented generation, and defines an integrated conceptual architecture. The proposed solution combines ESP32-based BLE beacons, smartphone sensors for orientation inference, voice interaction through TTS/STT, and an intelligence layer based on RAG to answer user queries with documentary support.

As a result, the thesis proposes a functional and economically defensible architecture for digital cultural mediation in the Peruvian context. Its main contribution lies in integrating proximity-based localization, mobile orientation, multimodal accessibility, and conversational intelligence within a low-cost BYOD approach that is scalable and compatible with the operational constraints of cultural institutions. In addition, the economic and financial assessment suggests favorable feasibility conditions for a progressive institutional deployment, supporting the technical and practical relevance of the proposal.

\noindent\textbf{Keywords:} smart museums, BLE, accessibility, RAG, digital cultural mediation
